%!TeX root=../main.tex

\chapter{Conclusioni}\label{ch:conclusioni}

\section{Sintesi del lavoro svolto}\label{sec:sintesi-del-lavoro-svolto}
Il progetto di tesi ha portato alla progettazione, implementazione e distribuzione di \textbf{Swift Contest}, una piattaforma software completa e funzionante per la gestione di concorsi, che raggiunge con successo tutti gli obiettivi prefissati in fase di analisi.
È stato dimostrato come, attraverso l'integrazione di un ecosistema tecnologico moderno e accessibile, sia possibile costruire un sistema complesso, sicuro e scalabile, capace di risolvere le inefficienze e le criticità dei processi di gestione tradizionali.

Il lavoro svolto, tuttavia, non si esaurisce con la presente tesi: esso costituisce una base di partenza per future iterazioni e, potenzialmente, per applicazioni in contesti reali, aprendo prospettive interessanti nel campo delle piattaforme collaborative e dei sistemi di votazione digitale.

\section{Risultati ottenuti}\label{sec:risultati-ottenuti}
I principali risultati ottenuti possono essere riassunti come segue:
\begin{itemize}
    \item \textbf{Realizzazione di un'applicazione cross-platform unificata}: utilizzando Flutter, è stato sviluppato un unico client che offre un'esperienza utente coerente e funzionale sia su piattaforme web che mobile (Android), validando l'efficacia di un approccio a codebase singola per ridurre i tempi di sviluppo e la complessità di manutenzione.
    \item \textbf{Implementazione di un'architettura backend sicura e robusta}: è stato costruito un backend basato sul Principio del Minimo Privilegio, dove l'accesso ai dati è interamente mediato da un'API controllata, proteggendo l'integrità del sistema.
    \item \textbf{Creazione di un sistema di votazione affidabile e trasparente}: il meccanismo di \emph{snapshot} delle sessioni di voto garantisce l'immutabilità dei dati e l'auditabilità del processo, risolvendo una delle principali criticità dei sistemi di valutazione manuali.
    \item \textbf{Validazione di un'infrastruttura efficiente e sostenibile}: l'architettura basata su servizi cloud (Supabase, Vercel) e l'uso mirato di risorse onerose come il \emph{realtime} dimostrano la fattibilità di costruire soluzioni complesse mantenendo i costi operativi contenuti, in linea con i piani gratuiti o a basso costo offerti dai provider.
\end{itemize}
Il prototipo realizzato non si limita a una dimostrazione tecnica, ma si configura come un'applicazione integrata e pronta per una fase di validazione sul campo, dimostrando la sua idoneità per contesti reali di piccola scala.

\section{Conoscenze acquisite}\label{sec:conoscenze-acquisite}
Durante lo sviluppo del progetto Swift Contest, ho avuto l'opportunità di approfondire e consolidare una vasta gamma di conoscenze e competenze, sia a livello tecnico che metodologico.
In particolare, ho acquisito una solida padronanza di:
\begin{itemize}
    \item \textbf{Sviluppo Frontend con Flutter e Dart}: Ho imparato a progettare e implementare interfacce utente complesse e reattive, gestendo lo stato dell'applicazione tramite il pattern BLoC~\cite{bloc_docs} e ottimizzando l'esperienza utente su diverse piattaforme.
    Ho approfondito l'uso di Dart Extensions per estendere le funzionalità di classi esistenti e la gestione della navigazione con AutoRoute~\cite{auto_route_package}.
    \item \textbf{Architetture Backend-as-a-Service (BaaS) con Supabase}: Ho compreso i vantaggi e le sfide di un'architettura BaaS, imparando a configurare e gestire un database PostgreSQL, a implementare logiche di business complesse tramite Funzioni RPC (PL/pgSQL) ed Edge Functions (TypeScript), e a sfruttare servizi come l'autenticazione, lo storage e il realtime.
    \item \textbf{Principi di Sicurezza del Software}: Ho applicato concretamente il Principio del Minimo Privilegio, progettando un accesso ai dati mediato da API e implementando meccanismi di autorizzazione robusti, inclusa la gestione dell'autenticazione anonima.
    \item \textbf{Design Pattern e Best Practice}: Ho utilizzato e compreso l'importanza di pattern come MVVM, BLoC, Singleton e la gestione centralizzata degli errori tramite il tipo Either, che contribuiscono a creare codice manutenibile e testabile.
    \item \textbf{Processi di Distribuzione e CI/CD}: Ho acquisito esperienza nella configurazione di pipeline di Continuous Integration/Continuous Deployment per applicazioni web e nella gestione sicura della distribuzione di APK Android tramite GitHub.
    \item \textbf{Modellazione del Dominio e Design del Database}: Ho imparato a tradurre requisiti complessi in un modello dati robusto, implementando strategie di integrità referenziale, cancellazione e snapshot dei dati per garantire consistenza e auditabilità.
\end{itemize}
Questa esperienza mi ha fornito una visione completa del ciclo di vita dello sviluppo software, dalla concezione all'implementazione e distribuzione.

\section{Valutazione personale}\label{sec:valutazione-personale}
La realizzazione di Swift Contest ha rappresentato un percorso di crescita significativo, non solo dal punto di vista tecnico, ma anche personale e progettuale.
Affrontare la sfida di costruire un sistema completo da zero mi ha permesso di apprezzare l'importanza di un'\textbf{architettura ben ponderata} fin dalle prime fasi.
In questo contesto, ho deliberatamente scelto di non integrare funzionalità legate all'intelligenza artificiale, pur riconoscendone l'enorme importanza nel panorama tecnologico attuale.
Per me era cruciale padroneggiare i \textbf{principi fondamentali} di un software ``tradizionale'' prima di aggiungere strati di complessità legati all'AI\@.

Se dovessi ripensare al progetto oggi, con l'esperienza maturata, dedicherei molta più attenzione alla \textbf{fase di progettazione iniziale}.
All'inizio del percorso, non possedevo la base di conoscenze che ho oggi sulle strategie moderne per affrontare problemi complessi come la gestione dello stato, la sicurezza delle API o la distribuzione.
Ho imparato sul campo che per progettare un sistema al meglio è fondamentale conoscere a fondo le alternative e i \textbf{pattern consolidati}.

Per questo motivo, sono estremamente soddisfatto del lavoro svolto, non solo per il risultato finale, ma soprattutto per il \textbf{valore formativo del processo}.
Questo progetto mi ha insegnato a \textbf{valutare i trade-off} e a scegliere gli strumenti giusti per il problema giusto.
Le competenze acquisite, in particolare nell'ambito dello sviluppo mobile e delle architetture backend, sono ora un bagaglio che posso trasferire con sicurezza nello sviluppo di qualsiasi software, sia che si tratti di utilizzare framework cross-platform alternativi come \textbf{Jetpack Compose Multiplatform}~\cite{compose_multiplatform_docs} e \textbf{Kotlin Multiplatform}~\cite{kotlin_multiplatform_docs}, sia che si tratti di approfondire lo sviluppo nativo, in particolare Android con \textbf{Jetpack Compose}~\cite{jetpack_compose_docs}.